\section*{Chimie (7 points)}
Les réactions acide-base et les réactions d'oxydo-réduction peuvent évoluer dans deux sens différents conduisant à la formation d'autres corps et/ou d'autres solutions. Ainsi, l'acide méthanoïque connu pour son odeur piquante donne par estérification un ester à l'odeur fruitée évoquant l'ananas, alors que le fonctionnement de la pile (Cuivre / Aluminium) conduit à un dépôt de cuivre.\\
Cet exercice vise :

\begin{itemize}
  \item l'étude d'une solution aqueuse d'acide méthanoïque ;
  \item l'étude d'une réaction d'estérification ;
  \item l'étude du fonctionnement de la pile (Cuivre/Aluminium).
\end{itemize}

\section*{Les parties 1 et 2 sont indépendantes}
\section*{Partie 1: Réactivité de l'acide méthanoïque}
On dispose d'une solution aqueuse ( $S_{A}$ ) d'acide méthanoïque $H_{C O O H_{(a q)}}$ de volume $V$ et de concentration molaire $C_{A}=2.10^{-2} ~mol . L^{-1}$. Le $p H$ de cette solution à $25^{\circ} C$ est $p H=2,75$.

\begin{enumerate}
  \item Écrire l'équation de la réaction entre l'acide méthanoïque et l'eau.
  \item Dresser le tableau d'avancement de cette réaction en utilisant les grandeurs $C_{A}, V$ et l'avancement final $x_{f}$ de la réaction.
  \item Calculer la valeur du taux d'avancement final $\tau$ de la réaction. Déduire.
  \item Calculer la valeur de $Q_{r, e ́ q}$, quotient de réaction à l'état d'équilibre.
  \item Vérifier que la valeur de $p K_{A}$ du couple $\left(H C O O H_{(a q)} / H C O O_{(a q)}^{-}\right)$est $p K_{A}=3,76$.
  \item Pour vérifier la valeur de la concentration $C_{A}$, on dose le volume $V_{A}=10 m L$ de la solution ( $S_{A}$ ) par une solution aqueuse $\left(S_{B}\right)$ d'hydroxyde de sodium $\left(Na_{(a q)}^{+}+HO_{(a q)}^{-}\right)$de concentration molaire $C_{B}=10^{-2} ~mol . L^{-1}$. La figure ci-contre donne la courbe $p H=f\left(V_{B}\right)$ obtenue lors de ce dosage.\\
6.1. Écrire l'équation de réaction de dosage supposée totale.\\
6.2. En exploitant la courbe de dosage, retrouver la valeur de $C_{A}$.
  \item Pour synthétiser l'ester à l'arôme d'ananas, on fait réagir l'acide méthanoïque pur $HCOOH_{(\ell)}$ avec le butan-1-ol $C_{4} H_{9} OH_{(\ell)}$.\\
7.1. Écrire, en utilisant les formules semi-développées, l'équation de la réaction d'estérification. Nommer l'ester formé.\\
7.2. Donner les caractéristiques de cette réaction.\\
7.3. Proposer deux méthodes pour augmenter le rendement de la synthèse de cet ester.\\
\end{enumerate}
\begin{figure}[H]
    \centering
    \includegraphics[width=.7\textwidth]{f47096eb-5ae7-4c60-ab2f-f9822564c423-2_840_697_1516_1302}
    \caption{}
    \label{fig:}
\end{figure}



\section*{Partie 2: Pile (Cuivre/Aluminium)}
On dispose d'une solution aqueuse ( $S_{A}$ ) d'acide méthanoïque $H_{C O O H_{(a q)}}$ de volume $V$ et de concentration molaire $C_{A}=2.10^{-2} ~mol . L^{-1}$. Le $p H$ de cette solution à $25^{\circ} C$ est $p H=2,75$.

\begin{enumerate}
  \item Écrire l'équation de la réaction entre l'acide méthanoïque et l'eau.
  \item Dresser le tableau d'avancement de cette réaction en utilisant les grandeurs $C_{A}, V$ et l'avancement final $x_{f}$ de la réaction.
  \item Calculer la valeur du taux d'avancement final $\tau$ de la réaction. Déduire.
  \item Calculer la valeur de $Q_{r, e ́ q}$, quotient de réaction à l'état d'équilibre.
  \item Vérifier que la valeur de $p K_{A}$ du couple $\left(H C O O H_{(a q)} / H C O O_{(a q)}^{-}\right)$est $p K_{A}=3,76$.
  \item Pour vérifier la valeur de la concentration $C_{A}$, on dose le volume $V_{A}=10 m L$ de la solution ( $S_{A}$ ) par une solution aqueuse $\left(S_{B}\right)$ d'hydroxyde de sodium $\left(Na_{(a q)}^{+}+HO_{(a q)}^{-}\right)$de concentration molaire $C_{B}=10^{-2} ~mol . L^{-1}$. La figure ci-contre donne la courbe $p H=f\left(V_{B}\right)$ obtenue lors de ce dosage.\\
6.1. Écrire l'équation de réaction de dosage supposée totale.\\
6.2. En exploitant la courbe de dosage, retrouver la valeur de $C_{A}$.
  \item Pour synthétiser l'ester à l'arôme d'ananas, on fait réagir l'acide méthanoïque pur $HCOOH_{(\ell)}$ avec le butan-1-ol $C_{4} H_{9} OH_{(\ell)}$.\\
7.1. Écrire, en utilisant les formules semi-développées, l'équation de la réaction d'estérification. Nommer l'ester formé.\\
7.2. Donner les caractéristiques de cette réaction.\\
7.3. Proposer deux méthodes pour augmenter le rendement de la synthèse de cet ester.\\
\end{enumerate}
\begin{figure}[H]
    \centering
    \includegraphics[width=.7\textwidth]{f47096eb-5ae7-4c60-ab2f-f9822564c423-2_840_697_1516_1302}
    \caption{}
    \label{fig:}
\end{figure}


